% ECONOMICS_PROBLEM: {"id": "E43-463", "title": "Secular stagnation and global saving–investment imbalance", "jel_code": "E43", "source_id": "OP-0018"}
% !TeX program = xelatex
\documentclass[11pt,a4paper]{article}
\usepackage[margin=23mm]{geometry}
\usepackage{fontspec}
\setmainfont{Latin Modern Roman}
\usepackage{amsmath,amssymb,mathtools}
\usepackage{xcolor,enumitem,booktabs,longtable,array,xurl,hyperref}
\definecolor{muted}{HTML}{53636C}
\hypersetup{colorlinks=true,urlcolor=blue,linkcolor=blue}
\setlength{\parindent}{0pt}
\setlength{\parskip}{5pt}
\newcommand{\R}{\mathbb R}
\newcommand{\E}{\mathbb E}
\newcommand{\N}{\mathbb N}
\newcommand{\ind}{\mathbf 1}
\newcommand{\Var}{\operatorname{Var}}
\newcommand{\Cov}{\operatorname{Cov}}
\newcommand{\supp}{\operatorname{supp}}
\DeclareMathOperator*{\argmin}{arg\,min}
\DeclareMathOperator*{\argmax}{arg\,max}
\newcommand{\entrymeta}[3]{{\sffamily\small\color{muted}\textbf{\texttt{#1}}\quad|\quad #2\par #3\par}}
\newcommand{\Problemref}[1]{\texttt{#1}}
\newcommand{\JELref}[2]{\texttt{#2} (JEL \texttt{#1})}
\begin{document}
\section*{Secular stagnation and global saving–investment imbalance}
\textbf{Problem ID:} \texttt{E43-463}\quad \textbf{JEL:} \texttt{E43}\par
% BEGIN ECONOMICS STATEMENT
\label{problem:OP-0018}
\label{atlas:prob:A8}

\noindent\textbf{Original atlas ID:} A8. \textbf{Formulation:} editorial specification of a research family.

\subsection*{Definitions and assumptions}

In a declared full-employment overlapping-generations or asset-market model, let $S(r;\theta)$ and $I(r;\theta)$ be desired saving and investment at real rate $r$, where $\theta$ includes demography, income distribution, technology, asset-safety premia, and international positions. Define excess desired saving $E(r;\theta)=S(r;\theta)-I(r;\theta)$ after correctly including net foreign asset flows or aggregating the world economy. A natural equilibrium rate $r^*(\theta)$ is a specified root $E(r;\theta)=0$. A benchmark rate $r_0$ makes ``normal'' precise. In a log-rate monetary approximation with a zero nominal lower bound and fixed expected inflation $\bar\pi$, the attainable real-rate lower bound is $r_{\min}=-\bar\pi$; other regimes require their own rate constraint.

\subsection*{Formal research question}

Identify which primitive changes account for a persistent fall in $r^*(\theta)$ or for $E(r_0;\theta)>0$, and whether the evidence supports $r^*(\theta)<r_{\min}$. Under differentiability, local uniqueness, and $E_r(r^*;\theta)\ne0$, the comparative-static target is
\[
\frac{\partial r^*}{\partial\theta_j}=-\frac{E_{\theta_j}(r^*;\theta)}{E_r(r^*;\theta)}.
\]
For a stated path $\theta(u)$, $u\in[0,1]$, decompose the rate change by integrating these derivatives times $\theta_j'(u)$. Characterize identified contributions using household saving by age, investment, yields, and cross-country assets. Report dependence on the counterfactual path and equilibrium selection; a numerical calibration alone does not identify the causal contributions.

\subsection*{Interpretation and scope}

Desired saving exceeding desired investment at a counterfactual rate is compatible with equilibrium accounting: realized saving and investment obey the relevant national-account identities, including external balances. The atlas must therefore be read as an ex-ante excess-demand question, not a violation of accounting. A sign claim about the required rate adjustment additionally needs the sign of $E_r$ and a suitable root. Safe-asset demand and productive-capital investment are different assets and must not be collapsed without an explicit intermediation mechanism. Persistent demand shortfalls at a lower bound require the nominal adjustment and expectation assumptions that prevent another route to full employment. Nonlinear mechanism decompositions can depend on the order or path of primitive changes.

\subsection*{Source audit and status}

Hansen (1939) and Summers (2014) are verified historical and policy arguments. Eggertsson--Mehrotra (2014) is verified as NBER Working Paper 20574, not a two-author published journal paper. Its later quantitative extension includes Jacob A. Robbins and was published in 2019; it is added as [S4]. The publisher and archival URLs distinguish these versions. Existing models already demonstrate secular-stagnation mechanisms under assumptions. The editorial remaining agenda is identification and discrimination of such mechanisms, with correct equilibrium accounting and regime dependence. Neither ``normal rates'' nor a universal persistent saving surplus supplies a fully specified open theorem in the source atlas.

\subsection*{References}

\begin{enumerate}

\item[S1] Alvin H. Hansen (1939). \emph{Economic Progress and Declining Population Growth}. American Economic Review 29(1), 1--15. \url{https://www.jstor.org/stable/1806983}. \textit{Locator:} presidential address. \textit{Establishes:} A historical stagnation argument involving population and investment. \textit{Does not establish:} a contemporary identified quantitative decomposition.

\item[S2] Lawrence H. Summers (2014). \emph{U.S. Economic Prospects: Secular Stagnation, Hysteresis, and the Zero Lower Bound}. Business Economics 49(2), 65--73. \url{https://link.springer.com/article/10.1057/be.2014.13}. \textit{Locator:} abstract; NABE address. \textit{Establishes:} A secular-stagnation diagnosis and policy discussion. \textit{Does not establish:} a universal saving--investment imbalance theorem.

\item[S3] Gauti B. Eggertsson and Neil R. Mehrotra (2014). \emph{A Model of Secular Stagnation}. NBER Working Paper 20574. \url{https://www.nber.org/papers/w20574}. \textit{Locator:} abstract; overlapping-generations model. \textit{Establishes:} A model permitting persistent low natural rates. \textit{Does not establish:} the later three-author published quantitative model.

\item[S4] Gauti B. Eggertsson, Neil R. Mehrotra, and Jacob A. Robbins (2019). \emph{A Model of Secular Stagnation: Theory and Quantitative Evaluation}. American Economic Journal: Macroeconomics 11(1), 1--48. \url{https://www.aeaweb.org/articles?id=10.1257/mac.20170367}. \textit{Locator:} abstract; 56-period quantitative model. \textit{Establishes:} A published quantitative extension and decomposition. \textit{Does not establish:} that its calibration identifies every proposed stagnation channel.

\end{enumerate}

\subsection*{Common conventions: Source mathematical conventions}

Unless an entry states otherwise, agent and firm sets are finite. State and action spaces are measurable, strategies and outcome maps are measurable, probability laws are specified, and expectations exist for the targets under discussion. Infinite-horizon discount factors lie in $(0,1)$ and discounted objectives are finite. Norms, tolerances and complexity input sizes must be specified for computational comparisons. Constants or primitives that appear locally are inputs, not empirical facts asserted by this catalogue.

\textbf{Identification.} A statistical model $m\in\mathcal M$ implies an observable law $P_m$ and target $\theta(m)$. At law $P$, its identified set is
\[
\Theta(P;\mathcal M)=\{\theta(m):m\in\mathcal M,\ P_m=P\}.
\]
Point identification holds when this set is a singleton, or equivalently when $P_{m_1}=P_{m_2}$ implies $\theta(m_1)=\theta(m_2)$ for all compatible models. Sharp bounds mean bounds attained, or approached when the boundary is not attained, by compatible models. Writing this set is a definition, not a solution: the substantive task is to establish informative restrictions and derive or estimate the set. An unrestricted-confounding model may give only trivial bounds.

\textbf{Inference.} Identification, estimability and valid uncertainty quantification are separate questions. Fix $\alpha\in(0,1)$. A confidence procedure $C_n$ over a declared law class $\mathcal P$ has asymptotic uniform coverage at level $1-\alpha$ when
\[
\liminf_{n\to\infty}\inf_{P\in\mathcal P}\Pr_P\{\theta(P)\in C_n\}\geq1-\alpha.
\]
For set-identified targets the coverage event must be defined separately, for example coverage of the full identified set. If an entry uses identification-indexed models instead of uniquely indexed laws, the same coverage convention is applied over those models. Pointwise limits do not establish uniform validity, and asymptotic validity does not guarantee finite-sample precision. Sample sizes, dependence structures and weak-identification sequences are part of each application.

\textbf{Counterfactuals.} $Y_i(a)$ is a potential outcome under a specified intervention, including its timing, implementation and versions. It is not $Y_i$ conditional on voluntarily choosing $a$. A full intervention vector permits interference; shorthand scalar treatment notation does not silently impose no interference. Assignment, selection, attrition, anticipation and measurement must be specified. A claim about an instrument's local effect is not automatically a population policy effect.

\textbf{Economic equilibrium.} A counterfactual model includes best responses, constraints, feasibility, market clearing, state transitions and expectations consistent with the stated information. When several equilibria exist, either a declared selector is an input or the target is set-valued. Comparisons across policy regimes must use the stated selection convention. Reduced-form relationships are not presumed invariant after an equilibrium-changing intervention.

\textbf{Optimization and welfare.} Optimization is over the stated feasible set. If attainment is not established, $\sup$ or $\inf$ and $\varepsilon$-optimal choices replace an asserted maximizer or minimizer. For example, with value $V=\sup_{a\in\mathcal A}W(a)<\infty$, an $\varepsilon$-optimal set is $\{a\in\mathcal A:W(a)\geq V-\varepsilon\}$ for $\varepsilon>0$. Benefits, costs, risk and health or environmental outcomes can be aggregated only using declared units and conversion weights. Interpersonal utility weights, the treatment of fiscal transfers and foreign profits, discounting, and the behavioral welfare criterion are explicit inputs. A comparative-static derivative additionally requires existence and appropriate differentiability of the selected counterfactual.

\textbf{Sources.} Local labels [S1], [S2], and so forth restart in every question. Locators identify the relevant abstract, model, section or result at the level actually checked; a bibliographic or abstract-level verification is not represented as inspection of an entire proof. Working papers are distinguished from later publications and changing versions. Source summaries are paraphrased. All URL checks and editorial formulations refer to the audit date stated above.
% END ECONOMICS STATEMENT
\end{document}
